\documentclass{article}
 \usepackage{amsmath}
 \usepackage{amsthm}
 \usepackage{amssymb}
 \usepackage{pgfplots}
 \usepackage[utf8]{inputenc}
 \usepackage{amsfonts}
 \usepackage[margin=2.5cm]{geometry}
 \usepackage{graphicx}
 \usepackage[export]{adjustbox}
 \usepackage{fancyhdr}
 \usepackage[portuguese]{babel}
 \usepackage{hyperref}
 \usepackage{lastpage}
 \usepackage{mathtools}
 \setcounter{section}{-1}

 \pagestyle{fancy}
 \fancyhf{}

 \pgfplotsset{compat = 1.18}

 \hypersetup{
     colorlinks,
     citecolor=black,
     filecolor=black,
     linkcolor=black,
     urlcolor=black
 }
 \newtheorem*{def*}{\underline{Defini\c c\~ao}}
 \newtheorem*{theorem*}{\underline{Teorema}}
 \newtheorem*{lemma*}{\underline{Lema}}
 \newtheorem*{prop*}{\underline{Proposi\c c\~ao}}
 \newtheorem{example}{\underline{Exemplo}}
 \newtheorem*{proof*}{\underline{Prova}}
 \renewcommand\qedsymbol{$\blacksquare$}
 \newcommand{\Lin}[1]{Lin_{\mathbb{K}}({#1})}

 \rfoot{P\'agina \thepage \hspace{1pt} de \pageref{LastPage}}

 \begin{document}
 \begin{figure}[ht]
  \minipage{0.76\textwidth}
    \includegraphics[width=4cm]{../icmc.png}
    \hspace{7cm}
    \includegraphics[height=4.9cm,width=4cm]{../brasao_usp_cor.jpg}
  \endminipage  
\end{figure}

\begin{center}
  \vspace{1cm}
  \LARGE
  UNIVERSIDADE DE S\~AO PAULO

  \vspace{1.3cm}
  \LARGE
  INSTITUTO DE CI\^ENCIAS MATEM\'ATICAS E COMPUTACIONAIS - ICMC

  \vspace{1.7cm}
  \Large
  \textbf{Notas de Espaços Métricos}

  \vspace{1.3cm}
  \large
  \textbf{Renan Wenzel - 11169472}

  \vspace{1.3cm}
  \large
  \textbf{Professora - Thaís Jordão}

  \textbf{E-mail: tjordao@icmc.usp.br}

  \vspace{1.3cm}
  \today
\end{center}

 \newpage

 \tableofcontents
 \newpage
 \section{Informações (Possivelmente) Úteis}
 \subsection{Datas das Provas:}
\begin{itemize}
  \item[P1)] 31/08 - Peso 1;
  \item[P2)] 03/10 - Peso 2;
  \item[P3)] 31/10 - Peso 2;
  \item[P4)] 23/11 - Peso 3;
  \item[P5)] 14/12 - Peso 3.
\end{itemize}
\subsection{Bibliografia}
 \begin{itemize}
   \item LIMA, E. L.  ``Espaços Métricos'', Rio de Janeiro: Projeto Euclides, 2005.
   \item DOMINGUES, H. H. ``Espaços Métricos e Introdução à Topologia'', Atual Editora, 1982.
 \end{itemize}
\subsection*{Monitoria}
  A ser definido. 

 \newpage

\section{Aula 01 - 08/08/2023}
\subsection{Motivações}
\begin{itemize}
  \item Introdução ao Material do Curso.
\end{itemize}
\subsection*{O que é um espaço métrico?}
  Ao longo deste curso, trabalharemos com um conjunto M não-vazio.
\hypertarget{def_metric}{ \begin{def*}
   Uma função \(d:M\times M\rightarrow \mathbb{R}\) é dita ser uma métrica em M se: 
  \begin{itemize}
    \item[i)] \(d(x, y)\geq 0, x, y\in M\);
    \item[ii)] \(d(x, y) = 0 \Leftrightarrow x = y, x, y\in M\);
    \item[iii)] \(d(x, y) = d(y, x), x, y\in M\);
    \item[iv)] \(d(x, y)\leq d(x, z) + d(z, y), x, y, z\in M\).
  \end{itemize}
  Neste caso, o par \((M, d)\) é chamado espaço métrico.
 \end{def*}
}
\begin{example}
 \begin{itemize}
   
   \item[1)] \((\mathbb{R}, d)\), em que \(d:\mathbb{R}\times \mathbb{R}\rightarrow [0, \infty)\) é dado por 
    d(x, y) = \(|x-y|.\) É claro que, olhando para \(d(x,y),\) vale para quaisquer x, y reais que 
      \[
        d(x, y) = |x-y| = |-1(y-x)| = 1|y-x| = d(y, x).
      \]
    Assim, resta verificarmos os itens dois e quatro da \hyperlink{def_metric}{definição de métrica}. Para o item (ii), 
    \[
      |x-y| = 0 \Longleftrightarrow x = y.
    \]
    Com relação ao último item, observe que 
    \[
      |x+y|\leq |x| + |y|.
    \]
    De fato, como \(|a|\geq a\) para todo número real a, 
    \[
      |x+y|^{2} = (x+y)^{2} = x^{2} + 2xy +y^{2}\leq x^{2} + 2|x||y| + y^{2} = (|x| + |y|)^{2}.
    \]
    Logo, tomando a raíz dos dois lados, segue a afirmação:
    \[
        |x+y|\leq |x| + |y|.
    \]
    Com isso, temos 
      \[
        d(x, y) = |x-y| = |x-z+z-y| = |(x-z)-(y-z)|\leq |x-z| + |y-z|.
      \]
    Portanto, \(d(x, y)\leq d(x, z) + d(z, y)\), o que torna \((\mathbb{R}, d)\) um espaço métrico.

    \item[2)] Seja X um conjunto não-vazio. Definimos 
      \[
        d:X\times X\rightarrow [0, \infty)
      \]
      por 
        \[
          d(x, y) = \left\{\begin{array}{ll}
              0,\quad x = y\\
              1,\quad x\neq y.
            \end{array}\right.
        \]
      Esta métrica é conhecida como \textit{métrica discreta}. Verifiquemos as propriedades dela.

Com efeito, como a imagem dela pode ser apenas 0 ou 1, o item 1 é trivial. Por definição, a métrica vale 0 se,
e somente se, x e y são iguais, tal que o item (ii) está feito. O item (iii) segue automaticamente se x e y são iguais. Caso
eles sejam diferentes, temos \(d(x, y) = 1, d(y, x) = 1\), ou seja, o item (iii) é válido para todos os casos. Por fim, a desigualdade triangular fica como exercício.
 \end{itemize} 
\end{example}
\newpage

\section{Aula 02 - 10/08/2023}
\subsection{Motivações}
\begin{itemize}
  \item a
\end{itemize}

 \end{document}
